% !PRINTED 393
\weberpage{393}
Danach haben wir:

4. \gesperrt{Ist $\frac14 D \equiv 1\ (\mathrm{mod}\ 8)$, so ist jede ungerade Zahl Normenrest von 4 und von 8.}

\gesperrt{Ist $\frac14 D \equiv 0\ (\mathrm{mod}\ 8)$, so sind nur die Zahlen von der Form $4n+1$ Normenreste von 4 und die Zahlen der Form $8n+1$ Normenreste von 8.}

\gesperrt{Ist $\frac14 D \equiv 2\ (\mathrm{mod}\ 4)$, so sind alle ungeraden Zahlen Normenreste von 4.}

\gesperrt{Ist $\frac14 D \equiv 2$, oder $\equiv 6\ (\mathrm{mod}\ 8)$, so sind im ersten Fall die Zahlen $8n+1$, $8n-1$, im zweiten Fall die Zahlen $8n+1$, $8n+3$ Normenreste von 8.}

\gesperrt{Ist $\frac14 D \equiv 3, 4, 7\ (\mathrm{mod}\ 8)$, so sind alle und nur die Zahlen der Form $4n+1$ Normenreste von 4 und von 8.}

5. \gesperrt{Ist $p$ eine ungerade Primzahl, und ein durch $p$ nicht teilbares $a$ Normenrest von $p^k$ (für irgend einen positiven Exponenten $k$), so ist $a$ auch Normenrest von $p^{k+1}$.}

Zum Beweis sei

\[
N(\omega) = a + p^k b,
\]
\[
\omega_1 = \omega\,(1 + x\,p^k),
\]

worin $x$ rational. Folglich

\[
\begin{aligned}
N(\omega_1)&=(a+p^k b)(1+xp^k)^2,\\
&\equiv a+p^k(b+ax)\ (\mathrm{mod}\ p^{k+1}),
\end{aligned}
\]

und wenn $x$ aus der Kongruenz $b + a x \equiv 0\ (\mathrm{mod}\ p)$ bestimmt wird, so folgt

\[
N(\omega_1) \equiv a\ (\mathrm{mod}\ p^{k+1}).
\]

6. \gesperrt{Ist $D$ ungerade, so gilt dasselbe auch noch für $p = 2$.}

Denn in diesem Falle ist jedes ungerade $a$ Normenrest von 8. Man setze also

\[
N(\omega) = a + 2^k b, \qquad k \geqq 3,
\]
\[
\omega_1 = \omega \left(1 + 2^k x\,\frac{1+\sqrt{D}}{2}\right),
\]
\[
N(\omega_1) \equiv (a + 2^k b)(1 + 2^k x)\ (\mathrm{mod}\ 2^{k+1}),
\]

und wenn man $a x + b \equiv 0\ (\mathrm{mod}\ 2)$ annimmt, so folgt

\[
N(\omega_1) \equiv a\ (\mathrm{mod}\ 2^{k+1}).
\]

7. \gesperrt{Ist $D$ gerade und $a$ Normenrest von $2^k \geqq 8$, so ist $a$ auch Normenrest von $2^{k+1}$.}
